\section{Introduzione}
La prima classificazione dei segnali distingue i segnali tra analogici e digitali. 

I segnali \textbf{analogici} sono continui nel tempo (prevediamo l'andamento tra -inf e +inf) 
e possono assumere un numero infinito di valori (un esempio di segnale analofico è la voce
modifca l'aria, le particelle che compongono l'aria e si propaga nello spazio; un altro esempio 
i segnali radio e elettrocardiogramma), vengono trasemssi e sono riproducibili, il segnale 
analogico per eccellenza è il vinile, il solco riproduce perfettamente le frequenze.
Il loro problema è che sono sensibili al rumore, se c'è un disturbo, il segnale viene
alterato e non è più possibile recuperare il segnale originale, è necessario quindi un 
segnale più robusto (segnale digitale che posso ricostruire anche se c'è un disturbo, 
basta che il disturbo non sia troppo grande).
Quando io passo passo da un segnale analogico a uno digitale ho fatto il
\textbf{campionamento}, estraggo i campioni a istanti di tempo.
Quasi tutti i segnali che trasmettiamo sono digitali, ma i segnali analogici campioanti
ancora non sono digitali (sequenza di 0 e 1).
Il campionamento mi dice che è vero che perdi informazione, ma Shannon ci dice che è 
possibile ricostruire il segnale originale (quello che sta in mezzo e viene perso).

I segnali \textbf{digitali} sono discreti nel tempo e possono assumere solo un numero 
finito di valori.

\section{Concetti base sui segnali}

Per spigare come funziona il teorema di Nyquist-Shannon, è necessario introdurre alcuni
concetti base sui segnali.
\subsection{Segnali sinusoidali}

\textbf{Periodo}: Intervallo di tempo dopo cui il segnale si ripete, si misura in secondi.

\textbf{Frequenza}: Numero di periodi al secondo, si misura in Hertz (Hz), è l'inverso del periodo.

\textbf{Ampiezza}: Valore massimo del segnale, si misura in Volt (V).

\textbf{Fase}: Punto di partenza del segnale, si misura in gradi o radianti.

\subsection{Serie di Fourier}
Se io ho un segnale periodico, comunque sia fatto (esempio onda quadra periodica), 
si dimostra che tutti i segnali periodici con periodo T possono essere rappresentati 
come somma di segnali sinusoidali (seno e coseno) con frequenze multiple di 1/T, 
cioè con frequenze discrete.\\
Questa rappresentazione è chiamata \textbf{serie di Fourier} e ci dice che qualsiasi 
segnale periodico può essere scomposto in una serie di segnali sinusoidali, ognuno 
con una frequenza, ampiezza e fase specifica.

Di solito non trasmetto segnali periodici, ma serve per introdurre trasformata di fourier.

Lo \textbf{spettro di ampiezza} e lo \textbf{spettro di fase} di un segnale 
sono rispettivamente il modulo e la fase dei coeffienti di Fourier.\\
Questa rappresentazione nel dominio della frequenza è complementare alla rappresentazione
del dominio del tempo.\\
Se conosco lo spettro di ampiezza e di fase conosco il segnale.\\
Dimensionamento del canale: capire quanta banda mi serve per trasmettere un segnale.

Effetto Gibbs: quando approssimo un segnale con una serie di Fourier, ottengo 
un'oscillazione vicino ai punti di discontinuità, che non scompare aumentando 
il numero di termini della serie, ma si riduce a una zona sempre più piccola 
intorno alla discontinuità.

\subsection{Trasformata di Fourier}
La \textbf{trasformata di Fourier} estende il concetto della Serie di Fourier ai segnali
non periodici.
 In questo corso la trasformata di fourier si chiama spettro del segnale.

Propietà della trasformata: ...

Nessun sistema di trasmissione è perfetto, la priam cosa che fa è attenuare il segnale. 
Se tutte le componenti di Fourier sono attenuate allo stesso modo, il segnale è 
attenuato ma non distorto, se invece alcune componenti sono più attenuate di altre, 
il segnale è distorto.

Tutti i canali di trasmissione hanno una banda ed è caratteristico di ogni sistema.
Il teorema di Shannon lega il bitrate massimo alla banda.

Nel wifi si usa un canale con una banda piccola.

\subsection{Conversione Analogico Digitale}
Il sistema deve prima campionare quindi trasformare il segnale analogico in uno tempo discreto
estrando dei campioni, Nyquist dice la freqeunza minima.

Per trasformare il numero in un dato digitale applico la quantizzazione.

Teorema di Nyquist ci dice prendi il segnale, vedi le frequenze, la frequenza di campionamento
deve essere almeno il doppio della frequenza massima del segnale. Sotto questa soglia si genera
l'aliasing, una distorsione che rende impossibile ricostruire il segnale originale.

L'errore di quantizzazione è la differenza tra il segnale analogico originale e il 
segnale digitale quantizzato, è un errore che si introduce quando si approssima un 
segnale continuo con un numero finito di livelli discreti.

Risoluzione in bit: con n bit posso rappresentare $2^n$ livelli di quantizzazione, 
più bit ho, più livelli di quantizzazione posso rappresentare, minore è l'errore 
di quantizzazione, ma maggiore è la quantità di dati da trasmettere.

\subsection{Teorema del campionamento}
Quando campione devo campionare con una frequenza maggiore o uguale la frequezna di Nyquist
$f_c=2W$.

\subsection{Frequenze di Campionamento}
Per segnali audio, la frequenza di campionamento standard è di 44.1 kHz, 
che è più del doppio della frequenza massima udibile (20 kHz). 
Per segnali video, la frequenza di campionamento dipende dalla risoluzione 
e dalla frequenza dei fotogrammi, ma è generalmente molto più alta.

\subsection{Quantizzazione}
Una volta campionato, devo decidere quanti intervalli di quantizzazione usare, 
più intervalli uso, più preciso è il segnale digitale, ma più dati devo trasmettere.

L'attuale standard per il segnale audio è di 16 bit (cioè 65536 livelli).

Di quanto spazio (byte) abbiamo bisogno per memorizzare un'ora di musica?
Devo trasmettere ogni secondo 44100 campioni, ogni campione è rappresentato da 16 bit 
(2 byte) al secondo, i canali audio sono due (stereo), quindi 
$$44100 campioni/s * 2 byte/campione * 2 canali = 176400 byte/s$$
Byte=635MB da memorizzare

Il bitrate per trasmettere il segnale audio correttamente?
$$44100 campioni/s * 16 bit/campione * 2 canali = 1411200 bit/s = 1411.2 kbps$$

I byte vengono usati per memorizzare i dati, i bit al secondo vengono usati per indicare
la velocità di trasmissione dei dati.

\section{Companding Non Lineare}
Il companding è una tecnica di quantizzazione non lineare che migliora il rapporto 
segnale-rumore (SNR) per segnali di bassa ampiezza.

\section{PCM Pulse Code Modulation}
Si prende il segnale registrato dal microfono di un telefono e si dice che la banda è 4 kHz,
non è vero la banda della voce è arriva pure a 100 kHz, ma per risparmiare banda si taglia a 4 kHz,
quindi 4 kHz è la banda della voce, la moltiplico per 2 (teorema di Nyquist) e ottengo 8 kHz, 
quindi campiono a 8 kHz (8000 campioni/s), quindi un campione ogni 125 microsecondi, ogni campione 
lo quantizzo e lo codifico in 8 bit (256 livelli di quantizzazione), quindi il bitrate è di 64 kbps 
$(8000 campioni/s * 8 bit/campione)$. 

Devo filtrare il segnale prima di campionarlo, altrimenti rischio di avere aliasing, 
quindi uso un filtro a 3.4 kHz e poi campiono, estraggo un campione, questo campione viene
inviato in un quantizzatore che trasforma il segnale in un numero digitale a 8 bit.

Questo è il processo di \textbf{PCM (Pulse Code Modulation)}, è un processo di conversione 
analogico-digitale, vale per qualsiasi segnale, non solo per la voce, ma anche per il video, 
l'audio, ecc.\\

In ottica si chiama on off keying, è un processo di modulazione digitale,
dove il segnale digitale viene modulato su un segnale portante, in questo caso la luce, 
accendendo e spegnendo la luce per rappresentare i bit 1 e 0.\\

Trasformo una sequenza di bit in un segnale che posso trasmettere, si chiama \textbf{modulazione}.

Devo permettere l'accesso multiplo, al canale di comunicazione di accedere a più utenti, 
ci sono vari dominii dipende da cosa devo fare. 

Utente 1 deve trasmettere 8 k campioni /s, quindi un campione ogni 125 microsecondi;\\
Nel frattempo (nei 125 microsendi) trasmetto altri utenti, in particolare ci trasmetto 32 utenti.
Time division multiplexing (TDM): divido il tempo in slot, ogni slot è di 125 microsecondi,
e ogni utente trasmette in uno slot, quindi ogni utente trasmette ogni 4 ms $(32 slot * 125 microsecondi/slot)$,
quindi ogni utente trasmette 250 campioni/s, quindi un bitrate di 2 kbps $(250 campioni/s * 8 bit/campione)$.

\subsection{Codifica di Sorgente}
Cerca di \textbf{ridurre il numero di bit necessari per rappresentare un segnale}, sfruttando la
ridondanza del segnale, ad esempio usando la codifica di Huffman, che assegna codici più 
brevi ai simboli più frequenti.\\

Ci sono due tipi: quelli con la perdita di informazione (lossy) e quelli senza perdita 
di informazione (lossless).

Primo teorema di Shannon: l'entropia è 
$$H_s=E\{I_k\}=\sum_{k=1}^{L} p_k \log_2 \frac{1}{p_k}= \sum_{k=1}^{L} p_k I_k$$

\subsection{Codifica di Canale}
Aggiunge bit, ovvero aggiunge ridondanza, per permettere la rilevazione e la 
correzione degli errori, per proteggerloi dagli errori introdotti dal rumore o dalle 
interferenze durante la trasmissione. Il suo scopo è garantire l'integrità del messaggio.\\

Fa riferimento al secondo teorema di Shannon, che stabilisce un limite superiore alla 
capacità di un canale di comunicazione, indicando la massima velocità alla quale è 
possibile trasmettere informazioni con una probabilità di errore arbitrariamente bassa, 
dato un certo livello di rumore.\\

$$C = B \log_2 (1 + \frac{S}{N})$$
dove:
- $C$ è la capacità del canale in bit al secondo (bps)
- $B$ è la banda del canale in hertz (Hz)
- $S$ è la potenza del segnale in watt (W)
- $N$ è la potenza del rumore in watt (W)
Il secondo teorema di Shannon ci dice che se vogliamo trasmettere a una velocità 
inferiore alla capacità del canale, possiamo farlo con una probabilità di errore 
arbitrariamente bassa, utilizzando codici di canale appropriati. 
Se invece vogliamo trasmettere a una velocità superiore alla capacità del canale, 
non importa quanto sofisticati siano i codici di canale che utilizziamo, non saremo 
in grado di ridurre la probabilità di errore al di sotto di un certo livello.

Nel gergo comune diciamo che non abbiamo banda, in realtà non abbiamo potenza del 
segnale, o abbiamo troppo rumore, quindi la capacità del canale è bassa. 

Teorema di Shannon: è possibile trasmettere dati attraverso un canale di comunicazione
affetto da rumore AWGN (Additive White Gaussian Noise) con bitrate R<C.

\paragraph{Rilevamento e correzione degli errori}
le tecniche di rilevazione degli errori sono due principali:
- Automatic repeat request (ARQ): il mittente invia i dati ed anche un codice a rilevazione
d'errore, che sarà utilizzato in ricezione per individuare gli eventuali errori, ed in 
tal caso cheidere la ritrasmissione dei dati corrotti. In modi casi la richiesta è
implicita; il destinatatio invia un acknowledgemebt (ACK) di corretta ricezione dei dati, 
ed il ...
- Forward error correction (FEC): il trasmettitore aggiunge bit di ridondanza al 
messaggio originale, in modo che il ricevitore possa rilevare e correggere gli errori 
senza dover richiedere una ritrasmissione, ad esempio utilizzando codici di correzione 
degli errori come i codici di Hamming o i codici Reed-Solomon.

Come controllo che c'è errore e ridondanza? 
Uso la codifica di hamming, si usa nella codifica di sorgente, quando ho estratto i campioni,
li codifico in bit, e se c'è un errore, la distanza di hamming mi dice quanti bit sono 
diversi tra il messaggio originale e quello ricevuto, se la distanza è 0, non c'è errore, 
se è 1, c'è un errore, se è 2, ci sono due errori, ecc.\\

Il codice di Hamming è un codice di correzione degli errori che può rilevare e 
correggere errori singoli, e rilevare errori doppi.\\

Forward Error Correction (FEC) è una tecnica di codifica che consente al ricevitore 
di correggere gli errori senza dover richiedere una ritrasmissione, utilizzando bit 
di ridondanza aggiunti al messaggio originale.\\

Tutti i sistemi di comunicazione introducono un errore e tutti i sistemi devono usare il BER
(Bit Error Rate) per misurare la qualità del canale, che è il numero di bit errati diviso
per il numero totale di bit trasmessi.\\

Il FEc è il sistema di codifica del canale e aggiunge bit di ridondanza per permettere la 
rilevazione e la correzione degli errori, causati nel canale di comunicazione. 
Ci sono vari tipi di FEC, quella sistematica e quella non sistematica.\\

Codici a blocchi, sono una categoria di algoritmi per la correzione degli errori (FEC) che
operano su segmenti di dati di lunghezza fissa. Il trasmettitore aggiunge bit di controllo 
a un blocco di bit di dati. Il ricevitore usa quei bit per verificare se ci sono stati errori e, 
se possibile, correggerli.

Tipi di codici a blocchi:
- Codici di Hamming: possono correggere errori singoli e rilevare errori doppi.
Il limite di questi codici è che devo mandare tanti bit per correggere pochi errori.
- Codici di Reed-Solomon: possono correggere errori multipli e sono ampiamente 
utilizzati in applicazioni come CD, DVD e comunicazioni satellitari.
Invece di tasmettere la sequenza di bit, considero che i miei simboli sono coefficienti
di un polimonio di grado k-1. Se ho un polimonio di k-1 o lo descrivo con i coeffiecienti 
o con i punti, per esempio x=1 il polimonio passa per un determinato punto.
Lavora prendendo una sequenza di 239 bit, ne aggiunge 16 per trasmettere 255 bit, e 
con 16 bit di controllo, correggo 8 bit.
Si riferisce a polimoni di gradi 8.
- Codici LDPC (Low-Density Parity-Check): sono codici a blocchi che offrono prestazioni 
vicine al limite di Shannon e sono utilizzati in standard di comunicazione come Wi-Fi e 5G.


Controllo degli errori in BLE:

Cyclic Redundancy Check (CRC)
Codici Convoluzionali


\section{Codifica di Linea}
Ci occupiamo della modulazione, e si può chiamare anche codifica di linea, 
tutti i sistemi che utilizziamo dobbiamo permettere agli utenti di accedere al mezzo 
i metodi per permettere la condivisione della risorsa fisica si chiama multiplexing o 
multiple access. 

Una volta che tutti gli utenti hanno trasmesso nel canale di tramessione in ricezione 
facciamo le operazioni inverse, non ci occupiamo più della codifica di sorgente, 
solo richiami della codifica di canale. Adesso entriamo nella parte più importante del 
corso.\\

La codifica di linea deve emeterre i bit supportante fisica, devo mettere il segnale digitale
sulla supportante fisica. Esempio filo di rame la supportante fisica è la tensione.

Se il segnale che voglio trasmettere è in fibra ottica, allora devo trasformare il segnake
in un segnale luminoso, lavoriamo nell'infrarosso.
In ricezione devo fare la demodulazione da ottico a elettrico.

Vediamo come deve essere fatto il segnale per rappresentare la sequenza di 0 e 1, 
on-off keying o non return to zero (NRZ), se voglio trasmettere un 1 accendo la luce, se voglio 
trasmettere uno 0 spengo la luce, è un sistema di modulazione digitale, ma non è molto 
efficiente, perché se ho una lunga sequenza di 0, non trasmetto nulla, e se ho una 
lunga sequenza di 1, trasmetto sempre la luce, quindi non c'è sincronizzazione tra il 
trasmettitore e il ricevitore, non so quando inizia e quando finisce un bit, quindi 
non è molto efficiente.\\

Un altro sistema è il return to zero (RZ), se voglio trasmettere un 1 accendo la luce 
per metà del tempo, e poi spengo la luce, se voglio trasmettere uno 0 spengo la luce per
tutto il tempo, è un sistema di modulazione digitale, ma non è molto efficiente,
perché se ho una lunga sequenza di 0, non trasmetto nulla, e se ho una lunga sequenza di 
1, trasmetto sempre la luce, quindi non c'è sincronizzazione tra il trasmettitore e il
ricevitore, non so quando inizia e quando finisce un bit, quindi non è molto efficiente.\\

Interferenza intersimbolica (ISI): quando il segnale di un simbolo interferisce con il 
segnale di un simbolo adiacente, causando distorsione e rendendo difficile la corretta 
interpretazione dei simboli ricevuti.\\

Criterio di Nyquist: per evitare l'interferenza intersimbolica, la forma d'onda del 
segnale deve essere progettata in modo tale che il segnale di un simbolo sia zero 
nei punti in cui il segnale di un simbolo adiacente è massimo.\\
La condizione di nyquist nel dominio del tempo è che la forma d'onda del segnale deve
 essere zero nei punti in cui il segnale di un simbolo adiacente è massimo, mentre 
 nel dominio della frequenza è che la banda del segnale deve essere limitata a metà
 della frequenza di simbolo.\\

Impulsi di Nyquist: sono forme d'onda che soddisfano il criterio di Nyquist e
 sono utilizzate per trasmettere simboli digitali senza interferenza intersimbolica. 
 Un esempio comune di impulso di Nyquist è l'impulso di sinc, che ha una forma a 
 campana e soddisfa il criterio di Nyquist sia nel dominio del tempo che nel dominio 
 della frequenza.\\

\subsection{Trasmissioni in banda base e banda passante}
La trasmissione in banda base si riferisce alla trasmissione di segnali digitali senza
modulazione, mentre la trasmissione in banda passante si riferisce alla trasmissione di
segnali digitali modulati su una portante a frequenza più alta.\\
La trasmissione in banda base è più semplice e meno costosa, ma è limitata dalla
distanza e dalla qualità del canale, mentre la trasmissione in banda passante 
consente di trasmettere segnali digitali su distanze più lunghe e attraverso canali 
di qualità inferiore, ma è più complessa e costosa.\\
....
....
....
\section{Multiplexing}
Deve permettere a piú utenti di accedere al canale di comunicazione. Questa 
tecnica è chiamata multiplexing o multiple access. Quindi é la tecnica di 
trasmettere insieme più segnali su un unico canale di comunicazione, in modo
che più utenti possano condividere la stessa risorsa fisica.\\

Quasi sempre le operazioni di multiplexing non modificano i segnali, ma li
combinano in modo da poterli trasmettere insieme.\\

Quale tecnologia possiamo usare per multiplexing? 
\begin{itemize}
    \item \textbf{Time Division Multiplexing (TDM)}: divide il tempo in slot e assegna ogni 
    slot a un utente diverso, permettendo a più utenti di trasmettere in momenti diversi.
    E' applicata nella telefonia, dove ogni chiamata viene assegnata a uno slot di tempo specifico, 
    e nella trasmissione dati, dove i pacchetti di dati vengono trasmessi in slot di tempo specifici.
    \item \textbf{Frequency Division Multiplexing (FDM)}: divide la banda del canale in frequenze 
    e assegna ogni frequenza a un utente diverso, permettendo a più utenti di trasmettere 
    contemporaneamente su frequenze diverse. E' applicata nella trasmissione radio, dove ogni 
    stazione radio trasmette su una frequenza specifica, e nella televisione, dove ogni canale 
    televisivo trasmette su una frequenza specifica.
    \item \textbf{Code Division Multiplexing (CDM)}: assegna a ogni utente un codice unico e permette 
    a più utenti di trasmettere contemporaneamente sulla stessa frequenza, distinguendo i
    segnali in base ai codici. E' applicata nelle comunicazioni militari, dove ogni unità militare
    ha un codice unico per comunicare, e nelle reti cellulari, dove ogni utente ha un codice unico per 
    comunicare con la torre cellulare.
\end{itemize}

Non é detto che ogni utente é l'utente finale, potrebbe essere un flusso di dati, ad esempio un video, 
che viene trasmesso da un server a un client, e il server potrebbe essere considerato un utente del 
canale di comunicazione.\\

\subsection{Time Division Multiplexing - TDM}
A ogni slot assento a un utente diverso, e mettendo insieme gli slot ottengo
un frame.\\

Devo trasmettere con una frequenza di campionamento 2 volte la banda del segnale, per trasmettere 
con correttezza, devo trasemttere 8000 byte al secondo.
Tra un segnale e un altro c'è un intervallo di 125 microsecondi, quindi in quel tempo posso trasmettere 
altri utenti, se ho 32 utenti, ogni utente trasmette ogni 4 ms, quindi ogni utente trasmette 250 campioni/s, 
quindi un bitrate di 2 kbps.\\

Se io non devo trasmettere dati quella risorsa viene persa, se un utente pagava di piú gli davo due slot
invece di uno, ma se non trasmetteva nulla, quella risorsa veniva persa, quindi non era efficiente.\\
Anche nelle reti di accesso, l'accesso si basa sul tempo, perché é il modo piú semplice 
di mettere insieme i dati.\\

Hanno inventato il multiplexing statistico, che non assegna slot fissi agli utenti, ma assegna slot 
dinamicamente in base alla domanda di traffico, migliorando l'efficienza del canale di comunicazione.\\

Come Funziona?\\
I segnali vanno memeorizzati e poi affasciati insieme in un frame, e poi trasmessi, in ricezione devo 
fare l'operazione inversa, estrarre i segnali dai frame, e poi decodificarli.\\

Quando la centrale riceve i dati, deve capire a quale utente appartengono, quindi ogni utente ha un identificatore unico,
che viene incluso nei dati trasmessi, in modo che la centrale possa instradare i dati al destinatario corretto.\\

Quando ricevo dalla centrale telefonica, so quale é il mio, so il mio tempo, e butto tutto il resto. 
E'una tecnica molto economica, perché non devo trasmettere dati di controllo, ma è inefficiente, 
perché se un utente non trasmette nulla, quella risorsa viene persa.\\

Confronto con PCM: \\
Nella TDM si é creato la Digital Hierarchy, che é una gerarchia di livelli di multiplexing, 
dove ogni livello combina più flussi di dati in un flusso di dati a velocità più alta.\\
Il livello più basso è il livello 0, che combina 24 canali di voce a 64 kbps ciascuno, per un totale di 1.544 Mbps (T1).\\
Il livello successivo è il livello 1, che combina 4 flussi T1 per un totale di 6.312 Mbps (E1).\\
Il livello successivo è il livello 2, che combina 4 flussi E1 per un totale di 34.368 Mbps (T2).\\
Il livello successivo è il livello 3, che combina 4 flussi T2 per un totale di 139.264 Mbps (T3).\\
Il livello successivo è il livello 4, che combina 4 flussi T3 per un totale di 565.148 Mbps (T4).\\
Il livello più alto è il livello 5, che combina 4 flussi T4 per un totale di 2.260.592 Mbps (T5).\\
 min 37


\section{Orthogonal Frequency Division Multiplexing - OFDM}
Non è un multiplexing, non significa che stiamo dando accesso a più utenti, è una tecnica 
di modulazione del segnale, è così potente che ha rivoluzionato i sistemi di comunicazione, 
è usato nella fibra ottica, è la base dei sistemi wi fi, nelle reti cellulari (4G e 5G), nelle trasmissioni televisive digitali,
ADSL/VDSL in quest'ultimi però si chiama DMT (Discrete Multitone Modulation).\\

OFDM è una tecnica di modulazione che divide un canale di comunicazione in più sottocanali ortogonali,
ciascuno dei quali trasmette una porzione del segnale complessivo, permettendo di trasmettere dati a 
velocità più elevate e migliorando la resistenza alle interferenze e alla distorsione del canale.\\

Quando trasmetto attraverso canale wireless, spesso il canale a causa del multipath fading, 
attenua alcune frequenze più di altre, questo causa distorsione del segnale, ma se allungo il tempo simbolo, 
restringo la banda e quindi lo spettro lo riduco, quindi il canale è più piatto, e quindi è più facile trasmettere, ma se allungo troppo il tempo simbolo,
aumenta la probabilità di interferenza intersimbolica, quindi OFDM divide il canale in più sottocanali (diverso) ortogonali, 
ognuno con una banda più stretta, permettendo di trasmettere dati a velocità più elevate e migliorando la resistenza 
alle interferenze e alla distorsione del canale.\\

Quindi prendo la banda disponibile, per esempio per wi fi ho 20 kHz, e la divido in 64 sottocanali, ognuno con una banda di 
312.5 Hz, e ogni sottocanale trasmette una porzione del segnale complessivo, quindi se ho 64 sottocanali, posso trasmettere 64 
simboli contemporaneamente, e quindi aumentare la velocità di trasmissione.\\

Efficienza spettrale: $\frac {R(bit-rate)}{B(banda)}$\\
L'efficienza spettrale di OFDM è elevata, perché permette di trasmettere dati a velocità più elevate utilizzando la stessa banda, 
migliorando l'efficienza spettrale rispetto ad altre tecniche di modulazione.\\
La sinc si annulla per valori come 1, 2 ecc, quindi l'idea è che se devo trasmettere un altro segnale nello spettro lo modulo
per farlo passare per quei valori, in questo modo non interferisce con il segnale precedente, e quindi posso trasmettere più segnali 
contemporaneamente senza interferenza.\\

La distanza tra i due segnali in frequenza è di $\frac 1 T$. 

Se il mio sistema di trasmissione si comporta strano in una delle frequenze posso rendermene conto e posso facilmente equalizzare.\\

\subsection{Multi-carrier modulation}
Non è un multiplexing, non è neanche una frequency division multiplexing, è una tecnica di modulazione che utilizza più portanti per trasmettere dati,
dobbiamo trasformare un segnale con R bit rate in N sottocanali (flussi) attraverso un Serial to Parallel Converter, e poi ciascun flusso
è trasmesso ad una diversa sottofrequenza (\textbf{subcarrier}).
Due sottocanali sono ortogonali se la loro frequenza è multipla di $\frac 1 T$ (distano tra loro di $\frac 1 T$), dove T è il tempo simbolo, 
in questo modo non interferiscono tra loro.\\

Ci aggiungo a questa sottoportante la portante vera, per esempio se la portante è a 2.4 GHz, e la sottoportante è a 312.5 Hz, allora 
la frequenza del segnale trasmesso sarà 2.4003125 GHz, e così via per gli altri sottocanali.\\

Discrete Multitone Modulation (DMT) è una tecnica di modulazione multi-carrier utilizzata principalmente nelle comunicazioni 
su linee telefoniche, come ADSL e VDSL, per trasmettere dati a velocità elevate su distanze relativamente lunghe. 
DMT divide la banda disponibile in più sottocanali, ognuno dei quali trasmette una porzione del segnale complessivo, 
migliorando l'efficienza spettrale e la resistenza alle interferenze e alla distorsione del canale.\\

Allargo la durata del bit, quindi la durata del simbolo è $\frac N R$, ho ristretto lo spettro. La separazione di ogni sinc 
è fatta in modo da non avere interferenza tra i sottocanali.\\

Come faccio a filtrarle visto che sono sovrapposte? Uso il principio di ortogonalità, i due segnali sono ortogonali se l'integrale del loro 
prodotto nella durata del simbolo è uguale a zero (ovvero delta di Dirac).

Scelgo gli esponenziali complessi come funzioni di base, sono chiamati kernel di Fourier, e come frequenza scelgo $\frac m T$, 
dove m è un intero, in questo modo ottengo la trasformata di Fourier discreta (DFT), che mi permette di trasformare il segnale dal dominio del 
tempo al dominio della frequenza, e viceversa.\\

Quindi se prendo la frequenza $\frac 1 T$ ho un coseno, se prendo la frequenza $\frac 2 T$ ho un coseno con frequenza doppia, e così via, 
in questo modo ottengo i sottocanali ortogonali.\\

Poichè se faccio trasformata di una rect ottengo una sinc, se la rect è moltiplicata per un coseno, ottengo una sinc traslata.
\subsection{Fast Fourier Transform - FFT}
La Fast Fourier Transform (FFT) è un algoritmo efficiente per calcolare la trasformata di Fourier discreta (DFT) di un segnale,
che è una rappresentazione del segnale nel dominio della frequenza.\\

Il processamento del segnale è fatto a livello digitale, ci faccio la trasformata di fourier e poi lo mando al DAC (lo trasforma ad analogico) per trasmetterlo.
S/P e P/S, per trasformare il segnale da seriale a parallelo e viceversa.\\

Le serie / parallelo dividono il flusso di dati (che hanno un bit rate alto) in più flussi di dati (con bit rate più basso) 
che possono essere trasmessi su N sottocanali ortogonali, quindi i vari flussi divisi avranno un bit rate di $\frac R N$, e poi alla fine li ricombino in seriale.\\
Le sottoportante le scelgo con frequenza $\frac 1 T$ perchè è l'inverso del tempo di simbolo.

Tornato da parallelo a seriale aggiungo un prefisso ciclico. Poi nella trasmissione devo fare la modulazione, quindi moltiplico per una portante, e poi trasmetto.\\
In ricezione è uguale ma prima devo demodulare, quindi moltiplico per la portante, poi tolgo il prefisso ciclico, poi faccio S/P, poi la FFT per 
trasformare dal dominio del tempo al dominio della frequenza, e poi faccio il parallelo a seriale per ricombinare i flussi di dati.\\

\paragraph{Prefisso ciclico}
Sto trasmettendo il segnale come flusso di bit, a causa della distorsione, questi bit si possono sovrapporre, causando l'interferenza intersimbolica, 
OFDM ha pensato che invece d mettere uno spazio vuoto tra un simbolo e l'altro, metto una copia del simbolo precedente (una parte del simbolo), 
in questo modo se c'è interferenza intersimbolica, il simbolo precedente è uguale a quello che sta arrivando, quindi non c'è interferenza, 
e posso ricostruire il segnale originale.\\

Più lungo è il prefisso ciclico, più è probabile che riesca a eliminare l'interferenza intersimbolica, ma più è lungo il tempo di simbolo, quindi
più è bassa l'efficienza spettrale, quindi c'è un tradeoff tra la lunghezza del prefisso ciclico e l'efficienza spettrale. \\%daje

Questo mi serve nei canali wireless per disperdere il multipath fading, nella fibra ottica per la dispersione cromatica, 
e in generale per qualsiasi canale che introduce distorsione.\\

\subsection{Schema del trasmettitore OFDM}
Il segnale di partenza lo divido in più flussi con convertitore S/P, in modo che ogni flusso abbia un Bit rate che è il bit rate di partenza diviso 
per il numero di flussi, creo una costellazione, ovvero 4 simboli e ogni simbolo rappresenta 2 bit (se ho 16 simboli ogni simbolo rappresenta 4 bit), 
OFDM poi fa il FFT per trasformare dal dominio del tempo al dominio della frequenza, poi P/S, poi aggiungo il prefisso ciclico, poi DAC,  
poi moltiplico per la portante (parte reale usa coseno, parte immaginaria usa coseno), e poi trasmetto.\\

In ricezione, prima moltiplico per la portante, lo filtro quindi tolgo la portante (rimetto in banda base), poi ADC, poi tolgo il prefisso ciclico, 
poi S/P, poi faccio la FFT per trasformare dal dominio del tempo al dominio della frequenza, poi faccio il parallelo a seriale per ricombinare i flussi di dati, 
e poi decodifico i simboli per ottenere i bit.\\

\subsection{Prestaziomi dell'OFDM}
PRO:
\begin{itemize}
    \item Robustezza rispetto all'interferenza intercanale, all'interferenza intersimbolica e al multipath fading, 
    grazie alla divisione del canale in sottocanali ortogonali e all'uso del prefisso ciclico.
    \item Efficienza spettrale elevata, grazie alla possibilità di trasmettere dati a velocità più elevate utilizzando la stessa banda.
    \item Permette formati di modulazione avanzati
    \item Implementazione efficiente, grazie all'uso della FFT per la modulazione e demodulazione dei sottocanali.
    \item Bassa sensibilità agli errori di sincronizzazione, grazie alla divisione del canale in sottocanali ortogonali e all'uso del prefisso ciclico.
    \item Processamento elettronico del segnale
\end{itemize}
CONTRO: 
\begin{itemize}
    \item Sensibile all'effetto doppler
    \item Sensibile ai problemi di sincronizzazione in frequenza
    \item Altro peak-to-average power ratio (PAPR), che può causare problemi di amplificazione del segnale e ridurre l'efficienza energetica del sistema.
    \item Perdita di efficienza dovuta all'inserimento del ciclo prefisso, che riduce la quantità di dati trasmessi per unità di tempo.
\end{itemize}

\section{Complementary Cumulative Distribution Function - CCDF}
La Complementary Cumulative Distribution Function (CCDF) è una funzione che descrive la probabilità che una variabile casuale superi un certo valore, 
ed è spesso utilizzata per analizzare il Peak-to-Average Power Ratio (PAPR) nei sistemi di comunicazione, come OFDM, per valutare la probabilità che 
il PAPR superi un certo livello, aiutando a progettare sistemi più efficienti e robusti contro le distorsioni causate da picchi di potenza elevati.\\

Peak to Average Power Ratio (PAPR): misura per quanto tempo il segnale che devo trasmettere supera il valore medio. 
CCDF è la probabilità che il PAPR superi un certo livello, ad esempio se voglio sapere la probabilità che il PAPR superi 10 dB, guardo la CCDF a 10 dB.\\

Il fatto che la dinamica del mio segnale è così ambia è un problema, si possono creare effetti non lineare, perché la potenza è troppo elevata e la risposta 
della fibra ottica non è lineare.\\

Questa tecnica la usiamo nel wi fi, in ottica e nei sistemi mobili.

\section{Discrete Multi Tone Modulation - DMT}
E' l'OFDM in banda banse e viene utilizzato nell'ADSL.
Nel 1960 abbiamo avuto i primi sistemi in rame, funzionava con un sistema analagico e quelle
rete di chiamavano POTS (Plain Old Telephone Service), si chiamavano old proprio perchè ora 
non esistono più, erano reti di comunicazione analogiche, e quindi non erano digitali, e quindi non potevano trasmettere dati, ma solo voce.\\

Il segnale voce era trasformato in segnale elettrico e trasmesso attraverso il filo di rame, e poi in ricezione veniva trasformato di nuovo in segnale acustico.\\

I primi 4 kHz non li tocchiamo (sono il segnale voce), dividiamo in sottoportanti la banda restante.

ADSL (Asymmetric Digital Subscriber Line) è una tecnologia di comunicazione che utilizza le linee telefoniche in rame 
per trasmettere dati a velocità elevate, e DMT è la tecnica di modulazione utilizzata in ADSL per dividere la banda 
disponibile in più sottocanali, ognuno dei quali trasmette una porzione del segnale complessivo, migliorando l'efficienza spettrale e l
a resistenza alle interferenze e alla distorsione del canale.\\

Considero 512 sottoportanti, ognuna con una banda di 4.3125 kHz, e ogni sottoportante trasmette una porzione del segnale complessivo, 
quindi se ho 512 sottoportanti, posso trasmettere 512 simboli contemporaneamente, e quindi aumentare la velocità di trasmissione.\\

Applico a tutti questi i simboli la trasformata di fourier, e poi li trasmetto, in ricezione faccio l'operazione inversa, 
quindi faccio la trasformata di fourier inversa, per trasformare dal dominio della frequenza al dominio del tempo, 
e poi decodifico i simboli per ottenere i bit.\\

Bit loading: il sistema misura i dati in trasmissione se c'è rumore o interferenza (che chiamiamo diafonia), 
e diminuisco il numero di bit che trasmetto (uso Shannon). Più c'è rumore, meno bit trasmetto, e viceversa.\\

Poi applico il QAM (Quadrature Amplitude Modulation) per modulare i simboli, poi trasformata di Fourier, poi +CP e poi DAC.\\

VDSL2 è quella che uso se a casa mia non arriva la fibra. Il tratto di rame deve essere più corto possibile, 
più è lungo più c'è il rumore, quindi più è basso il bit rate, e più è corto più è alto il bit rate.\\

I sistemi in banda base sono quelli che viaggiano in rame. 